\chapter{Node.js}

Sappiamo ora che le pagine web possono essere generate al volo da programmi in
esecuzione sul server. E sappiamo già molto di JavaScript: finora, però, l'abbiamo
usato soltanto nell'ambiente del browser.

Nulla vieta di usarlo anche in un \textbf{ambiente server}. Esistono diversi
runtime che lo permettono --- \term{Node.js}, \term{Deno}, \term{Bun} --- e in
questo corso useremo il primo, di gran lunga il più diffuso.

\section{Che cos'è Node.js}

\dfn{Node.js}{
  \term{Node.js} è un \textbf{ambiente di esecuzione} (\textit{runtime})
  JavaScript, open source e multipiattaforma. Esegue il motore JavaScript
  \textbf{V8} di Google, lo stesso del progetto Chromium, e adotta un modello di
  I/O \textbf{basato sugli eventi, asincrono per impostazione predefinita e non
  bloccante}.}

Per capire perché quel modello di I/O sia così importante, conviene partire da
come funzionano gli altri ambienti.

\section{Il modello multi-thread}

In molti ambienti e linguaggi, \textbf{una richiesta è gestita da un
thread o processo dedicato}, prelevato da un pool preallocato.

Il problema è che, per servire una richiesta, un'applicazione web deve spesso
compiere operazioni di I/O \textbf{lente e bloccanti}:

\begin{itemize}
  \item accedere a un database;
  \item accedere a un'API esterna (per esempio con \code{fetch()});
  \item leggere un file dal filesystem locale;
  \item leggere il corpo della richiesta.
\end{itemize}

Di conseguenza, i thread che gestiscono le richieste passano \textbf{molto
tempo ad aspettare} che queste operazioni si completino.

\begin{figure}[H]
  \centering
  \begin{tikzpicture}[font=\Lato\scriptsize,
      act/.style={fill=primary!30!white, draw=primary!60!black, line width=0.5pt,
                  minimum height=0.45cm, anchor=west, inner sep=0pt},
      wait/.style={fill=neutral!14!white, draw=neutral!45!white, line width=0.5pt,
                   minimum height=0.45cm, anchor=west, inner sep=0pt},
      io/.style={fill=orange-gradient-2!60!white, draw=orange-gradient-6,
                 line width=0.5pt, minimum height=0.35cm, anchor=west,
                 inner sep=0pt}]

    % barra del thread
    \node[act,  minimum width=0.9cm] at (0,0) {};
    \node[wait, minimum width=2.0cm] at (0.9,0) {};
    \node[act,  minimum width=0.7cm] at (2.9,0) {};
    \node[wait, minimum width=2.4cm] at (3.6,0) {};
    \node[act,  minimum width=0.7cm] at (6.0,0) {};
    \node[wait, minimum width=1.8cm] at (6.7,0) {};
    \node[act,  minimum width=0.9cm] at (8.5,0) {};

    \node[font=\Lato\scriptsize, anchor=east, text=neutral!25!black]
      at (-0.2,0) {thread};

    % operazioni di I/O
    \node[io, minimum width=2.0cm] at (0.9,-0.75) {};
    \node[io, minimum width=2.4cm] at (3.6,-0.75) {};
    \node[io, minimum width=1.8cm] at (6.7,-0.75) {};
    \node[font=\Lato\scriptsize, anchor=east, text=neutral!25!black, align=right]
      at (-0.2,-0.75) {operazioni\\ di I/O};

    % etichette
    \node[font=\Lato\scriptsize, text=neutral!35!black] at (1.9,0.4) {attesa};
    \node[font=\Lato\scriptsize, text=neutral!35!black] at (4.8,0.4) {attesa};
    \node[font=\Lato\scriptsize, text=neutral!35!black] at (7.6,0.4) {attesa};

    \draw[-{Stealth[length=5pt]}, draw=neutral!45!black, line width=0.7pt]
      (0,-1.35) -- (9.6,-1.35);
    \node[font=\Lato\scriptsize, text=neutral!35!black, anchor=north]
      at (9.6,-1.42) {tempo};

    \node[font=\Lato\scriptsize, text=primary!70!black, anchor=south]
      at (0.45,0.25) {};
  \end{tikzpicture}
  \caption{In azzurro il tempo in cui il thread lavora davvero; in grigio quello
           in cui aspetta.}
  \label{fig:blocking-io}
\end{figure}

Questo modello ha due difetti:

\begin{itemize}
  \item è \textbf{inefficiente}: moltissimo tempo di thread è speso ad
        aspettare;
  \item ha \textbf{scalabilità limitata}: il pool di thread è preallocato a una
        dimensione fissa, e non si possono gestire più richieste concorrenti di
        quante ne consenta il pool.
\end{itemize}

\subsection{Una metafora: il ristorante}

Il docente propone un'analogia che vale la pena riportare per esteso, perché
rende il concetto immediato.

Un ristorante ha \textbf{10 tavoli}. La direzione assume \textbf{10 camerieri},
assegnandone uno a ciascun tavolo.

\begin{itemize}
  \item Quando arriva un cliente, viene fatto accomodare al primo tavolo libero
        e il suo cameriere gli porta il menu. Il cliente impiega dai 5 ai 10
        minuti per decidere: \textbf{in questo tempo il cameriere resta fermo}.
  \item Deciso l'ordine, il cameriere lo porta in cucina. Il cuoco impiega dai
        15 ai 30 minuti per preparare: \textbf{il cameriere resta fermo}.
  \item Quando il piatto è pronto, il cameriere lo porta al tavolo. Il cliente
        impiega dai 15 ai 30 minuti per mangiare: \textbf{il cameriere resta
        fermo}.
  \item Finito il pasto, il cameriere sparecchia e chiede il conto alla cassa,
        che impiega dai 2 ai 5 minuti: \textbf{il cameriere resta fermo}.
  \item Preparato il conto, il cameriere lo porta al cliente, che paga e se ne
        va. Solo allora quel cameriere può servire un nuovo cliente.
\end{itemize}

I camerieri sono i \textbf{thread}, i clienti sono le \textbf{richieste}. Al
massimo 10 richieste possono essere servite contemporaneamente, e le attese
(decidere, cucinare, mangiare, preparare il conto) sono le operazioni
\textbf{bloccanti}. Per servire una richiesta in più bisogna assumere un nuovo
cameriere.

\subsection{Il ristorante efficiente}

Proviamo un'organizzazione diversa: il ristorante assume \textbf{un solo
cameriere} per tutti e 10 i tavoli.

\begin{itemize}
  \item Quando arriva un cliente, il cameriere lo accoglie, lo fa accomodare e
        gli porta il menu. \textbf{Mentre il cliente sceglie, il cameriere è
        libero} di occuparsi degli altri.
  \item Quando il cliente ha deciso, il cameriere porta l'ordine in cucina.
        \textbf{Mentre il cuoco prepara, il cameriere torna in sala} ed è
        disponibile.
  \item Quando il piatto è pronto (e la campanella della cucina suona), il
        cameriere lo porta al tavolo. \textbf{Mentre il cliente mangia, il
        cameriere è di nuovo libero.}
\end{itemize}

Un solo cameriere serve dieci tavoli senza problemi, perché non passa mai il
tempo ad aspettare. La campanella della cucina è la chiave: è un \textbf{evento}
che segnala che un'operazione lenta è terminata.

\section{Il modello di Node.js}

Node.js è basato su una filosofia molto simile a quella del ristorante
efficiente.

\begin{itemize}
  \item Un programma Node.js gira in un \textbf{unico processo}: non serve
        creare nuovi thread per ogni richiesta.
  \item Il runtime fornisce nella libreria standard un insieme di
        \term{primitive di I/O asincrone} che impediscono al codice JavaScript
        di bloccarsi.
  \item La maggior parte delle librerie adotta il paradigma non bloccante
        ovunque possibile, con ampio uso del pattern delle \textbf{callback}.
  \item Le operazioni sincrone sono l'\textbf{eccezione}, non la norma.
\end{itemize}

\begin{figure}[H]
  \centering
  \begin{tikzpicture}[font=\Lato\small,
      n/.style={wtnode, minimum width=3.0cm, minimum height=1.0cm},
      q/.style={wtnodeplain, minimum width=2.4cm, minimum height=0.9cm},
      a/.style={-{Stealth[length=5pt,width=4pt]}, line width=0.85pt,
                draw=neutral!45!black},
      c/.style={font=\Lato\scriptsize, anchor=east, text=neutral!25!black}]

    \node[c] (c1) at (-0.4,0.9)  {Client 1};
    \node[c] (c2) at (-0.4,0.1)  {Client 2};
    \node[c] (c3) at (-0.4,-0.9) {Client \textit{n}};
    \node[font=\Lato\small, text=neutral!45!black] at (-0.75,-0.4) {\vdots};

    \node[q] (ev) at (1.9,0) {coda\\ degli eventi};
    \node[n, draw=green-gradient-6, fill=green-gradient-1!45!white]
      (el) at (5.6,0) {\textbf{Event Loop}\\\scriptsize (un solo thread)};
    \node[q] (wt) at (9.4,0.9) {worker\\ thread};
    \node[wtnodeplain, minimum width=2.4cm, minimum height=0.9cm]
      (io) at (9.4,-1.0) {database,\\ filesystem, rete};

    \draw[a] (c1.east) -- (ev.west |- c1);
    \draw[a] (c2.east) -- (ev.west |- c2);
    \draw[a] (c3.east) -- (ev.west |- c3);
    \draw[a] (ev) -- (el);
    \draw[a] (el.east) -- (wt.west);
    \draw[a] (wt.south) -- (io.north);
    \draw[a, draw=accent!70!black] (io.west) to[out=180, in=-60]
      node[wtlab, pos=0.5]{callback} (el.south east);
  \end{tikzpicture}
  \caption{L'event loop a thread singolo di Node.js.}
  \label{fig:event-loop}
\end{figure}

\info{Un thread solo, ma non per tutto}{
  \guillemotleft Node.js è a thread singolo\guillemotright\ è una semplificazione utile ma
  imprecisa. È il \textbf{nostro} codice JavaScript a girare in un solo thread:
  le operazioni di I/O vere e proprie sono delegate a un pool di worker interni
  al runtime. La conseguenza pratica è importante: un calcolo lungo e
  \textbf{sincrono} nel nostro codice (un ciclo pesante, una cifratura,
  l'elaborazione di un'immagine) blocca l'event loop e quindi
  \textbf{tutte} le richieste, non solo la propria. Node.js eccelle nei carichi
  dominati dall'I/O, molto meno in quelli dominati dalla CPU.}

\section{Installare Node.js}

Ci sono tre modi:

\begin{itemize}
  \item usare gli installer disponibili sul sito ufficiale;
  \item usare il proprio gestore di pacchetti preferito;
  \item usare un \term{gestore di versioni} come \code{nvm}, \code{nvs} o
        \code{nvm4w}, comodo per sperimentare con versioni diverse e passare
        rapidamente dall'una all'altra.
\end{itemize}

Nel corso si usa Node.js \textbf{20.9.0}; con qualunque versione supportata non
dovrebbero esserci problemi. Per verificare che tutto funzioni:

\begin{affianco}[0.45][c]
\begin{lstlisting}[style=shell]
node --version
\end{lstlisting}
\risultato
\begin{console}[Terminale]
v20.9.0
\end{console}
\end{affianco}

\begin{esempio}{Hello World in Node.js}
\begin{affianco}[0.45][c]
\begin{lstlisting}[style=js, name=hello-world.js]
"use strict";
console.log("Hello Node.js!");
\end{lstlisting}
\risultato
\begin{console}[Terminale]
$ node hello-world.js
Hello Node.js!
\end{console}
\end{affianco}
\end{esempio}

\subsection{Differenze rispetto all'ambiente browser}

\begin{center}
\renewcommand{\arraystretch}{1.4}
\begin{tabular}{@{}>{\raggedright\arraybackslash}p{6.3cm}@{\hspace{1.0cm}}>{\raggedright\arraybackslash}p{6.3cm}@{}}
  \textbf{Nel browser} & \textbf{In Node.js} \\
  JavaScript serve soprattutto a manipolare il DOM: ci sono \code{window},
  \code{document}, le API della piattaforma web come i cookie e
  \code{localStorage}.
  & Questi oggetti e queste API \textbf{non esistono}. In compenso c'è la
    libreria standard con i suoi moduli per l'accesso al filesystem e alla
    rete: si possono leggere e scrivere file, e mettersi in ascolto su un socket
    per ricevere richieste. \\
\end{tabular}
\end{center}

\section{npm: il gestore di pacchetti}

Uno dei principali motori del successo di Node.js è \term{npm}. A settembre 2022
erano disponibili oltre 2,1 milioni di pacchetti: il più grande repository per
un singolo linguaggio al mondo. Esiste un pacchetto per (quasi) tutto.

npm viene installato insieme a Node.js. Per verificarlo:

\begin{lstlisting}[style=shell]
npm --version
\end{lstlisting}

\subsection{Creare un pacchetto}

Per creare un pacchetto npm nella directory corrente si usa \code{npm init}.
Poiché in questo corso lavoreremo con i \textbf{moduli ES}, useremo la variante:

\begin{lstlisting}[style=shell]
npm init es6
\end{lstlisting}

Il comando chiede alcune informazioni sul progetto (nome, versione, repository,
licenza, autore, punto di ingresso, comando di test) e al termine crea un file
\code{package.json} nella stessa directory, che contiene tutte le informazioni
sul pacchetto, comprese le dipendenze e i comandi per eseguirlo. Si può pensarlo
come la versione npm del \code{pom.xml} di Maven.

\begin{affianco}[0.6][c]
\begin{lstlisting}[style=json, name=package.json]
{
  "name": "hello-web-technologies",
  "version": "0.0.1",
  "description": "Web Technologies' first npm package",
  "main": "app.js",
  "type": "module",
  "scripts": {
    "test": "echo \"Error: no test specified\" && exit 1"
  },
  "keywords": ["greetings"],
  "author": "Luigi Libero Lucio Starace",
  "license": "MIT"
}
\end{lstlisting}
\risultato
\begin{immagine}[Le tre voci da conoscere.]
\begin{tikzpicture}[
    v/.style={rounded corners=3pt, draw=primary!60!black, fill=primary!8!white,
      line width=0.7pt, font=\FiraCodeNL\scriptsize, anchor=west, minimum width=1.9cm},
    d/.style={mnota, anchor=west, text width=3.4cm}]
  \node[v] (m) at (0,0)     {"main"};
  \node[d] at (2.1,0)       {il file che esporta ciò che il pacchetto offre};
  \node[v] (t) at (0,-1.0)  {"type"};
  \node[d] at (2.1,-1.0)    {\code{"module"} abilita \code{import} ed \code{export}};
  \node[v] (s) at (0,-2.0)  {"scripts"};
  \node[d] at (2.1,-2.0)    {comandi eseguibili con \code{npm <nome>}};
\end{tikzpicture}
\end{immagine}
\end{affianco}

\begin{itemize}
  \item il \textbf{file principale} (\code{app.js} nell'esempio) viene eseguito
        quando il codice client importa il nostro pacchetto: di norma esporta le
        variabili e le funzioni pubbliche;
  \item la proprietà \code{scripts} permette di definire \textbf{attività} da
        riga di comando, eseguibili con \code{npm <nome>};
  \item \code{"type": "module"} è ciò che abilita i moduli ES, cioè la sintassi
        \code{import}/\code{export} vista nel capitolo 6.
\end{itemize}

Aggiungiamo un'attività \code{start}:

\begin{affianco}[0.55][c]
\begin{lstlisting}[style=json, name=package.json]
"scripts": {
  "start": "node app.js",
  "test": "echo \"Error: no test specified\" && exit 1"
},
\end{lstlisting}
\risultato
\begin{console}[Terminale]
$ npm start

> hello-web-technologies@0.0.1 start
> node app.js

Hello Web Technologies!
\end{console}
\end{affianco}

\subsection{Installare dipendenze}

Per installare una dipendenza si usa \code{npm install <nomepacchetto>}.
Proviamo con \code{cowsay}:

\begin{lstlisting}[style=shell]
npm install cowsay
\end{lstlisting}

Succedono due cose:

\begin{enumerate}
  \item npm tiene traccia della nuova dipendenza, aggiungendo una sezione
        \code{dependencies} al file \code{package.json} (a sinistra);
  \item le dipendenze vengono scaricate nella directory \code{node\_modules}
        (a destra). Si noti che npm ha scaricato non solo il pacchetto
        richiesto, ma anche le sue dipendenze (ricorsivamente).
\end{enumerate}

\begin{affianco}[0.4][c]
\begin{lstlisting}[style=json]
"dependencies": {
  "cowsay": "^1.5.0"
}
\end{lstlisting}
\risultato
\begin{immagine}[Una dipendenza dichiarata, molte scaricate:\\
                 anche quelle di \code{cowsay}, ricorsivamente.]
\begin{tikzpicture}[
    dir/.style={mcod, anchor=west, inner sep=1.5pt},
    sel/.style={dir, fill=accent!14!white, draw=accent!55!black, rounded corners=2pt},
    ramo/.style={draw=neutral!50!black, line width=0.5pt}]
  \node[dir] (r) at (0,0) {node\_modules/};
  \foreach \n/\st [count=\i] in {cowsay/sel, get-stdin/dir, string-width/dir,
      strip-final-newline/dir, yargs/dir, \ldots/dir} {
    \node[\st] (n\i) at (0.5,-\i*0.48) {\n};
    \draw[ramo] (0.25,-\i*0.48+0.38) |- (n\i.west);
  }
  \node[mnota, anchor=west] at ([xshift=6pt]n1.east) {quella richiesta};
  \node[mnota, anchor=west, text=neutral!45!black] at ([xshift=6pt]n3.east) {le sue dipendenze};
\end{tikzpicture}
\end{immagine}
\end{affianco}

\begin{affianco}[0.48][c]
\begin{lstlisting}[style=js, name=app.js]
import cowsay from 'cowsay';

console.log(cowsay.think({
  text: "Hello Web Tech!"
}));
\end{lstlisting}
\risultato
\begin{console}[Terminale]
 _________________
( Hello Web Tech! )
 -----------------
        o   ^__^
         o  (oo)\_______
            (__)\       )\/\
                ||----w |
                ||     ||
\end{console}
\end{affianco}

\warning{node\_modules non va versionato}{
  Quando si distribuisce un pacchetto npm basta fornire il \textbf{proprio}
  codice e il descrittore \code{package.json}: la directory
  \code{node\_modules} \textbf{non} va aggiunta al controllo di versione. Chi
  riceve il progetto esegue \code{npm install} e npm ricostruisce l'intero
  albero delle dipendenze a partire dal \code{package.json}. È il motivo per cui
  \code{node\_modules/} è la prima riga di quasi ogni \code{.gitignore}.}

\section{Una prima applicazione web con Node.js}

\begin{esempio}{Un server HTTP in venti righe}
\begin{affianco}[0.6][c]
\codicepiccolo
\begin{lstlisting}[style=js, name=app.js]
import http from 'http';

const PORT = 3000;

let server = http.createServer(function (request, response) {
  let course = "Web Technologies";

  response.writeHead(200, { "Content-Type": "text/html" });
  response.write(`<!DOCTYPE html>
    <html><body>
      <h1>Hello ${course}</h1>
      <p>La data corrente e' ${new Date().toDateString()}</p>
    </body></html>`);
  response.end();
}).listen(PORT);

console.log(`Web app in ascolto sulla porta ${PORT}`);
\end{lstlisting}
\risultato
\begin{browser}[localhost:3000]
  {\raggedright\Lora\bfseries\LARGE Hello Web Technologies\par}\vspace{6pt}
  La data corrente e' Mon Oct 05 2026
\end{browser}
\begin{immagine}[Ogni richiesta invoca la callback,\\ che scrive la risposta.]
\begin{tikzpicture}[
    att/.style={rounded corners=3pt, draw=primary!65!black, fill=primary!12!white,
      line width=0.7pt, font=\Lato\scriptsize\bfseries, minimum width=1.6cm,
      minimum height=0.5cm},
    req/.style={-{Stealth[length=5pt,width=4pt]}, line width=0.8pt, draw=primary!70!black},
    res/.style={-{Stealth[length=5pt,width=4pt]}, line width=0.8pt, draw=green-gradient-6},
    msg/.style={mcod, above}]
  \node[att] (b) at (0,0) {browser};
  \node[att] (s) at (4.3,0) {app.js};
  \draw[req] ([yshift=4pt]b.east) -- node[msg]{GET /} ([yshift=4pt]s.west);
  \draw[res] ([yshift=-4pt]s.west) -- node[mcod, below]{200 text/html} ([yshift=-4pt]b.east);
\end{tikzpicture}
\end{immagine}
\end{affianco}

Qui si vede bene la differenza rispetto al server statico del capitolo 2: non
c'è alcuna document root, non c'è alcun file HTML sul disco. Il documento viene
\textbf{costruito a ogni richiesta}, e la funzione passata a
\code{createServer} è una \textbf{callback}, invocata dall'event loop ogni volta
che arriva una richiesta.

Si riconoscono anche gli elementi del protocollo studiati nel primo capitolo:
il codice di stato \code{200}, l'header \code{Content-Type}, il corpo della
risposta.
\end{esempio}

\section{Strumenti di sviluppo}

\subsection{Debug in Visual Studio Code}

Il modo più semplice è abilitare la funzione \term{Auto Attach}, che collega un
debugger ai processi Node.js avviati dentro VS Code. Per abilitarla si preme
\code{CTRL+SHIFT+P} per aprire la palette dei comandi e si cerca
\guillemotleft Toggle Auto Attach\guillemotright; conviene poi selezionare la modalità
\textbf{Smart}, che si collega automaticamente solo agli script fuori da
\code{node\_modules}.

I punti di interruzione si aggiungono come di consueto, facendo clic accanto al
numero di riga. Quando l'esecuzione raggiunge un punto di interruzione si ferma,
e il pannello \textit{Run and Debug} permette di ispezionare le variabili,
osservare espressioni, esaminare la pila delle chiamate e controllare il flusso.

\begin{center}
\renewcommand{\arraystretch}{1.4}
\begin{tabular}{@{}l@{\hspace{1.0cm}}>{\raggedright\arraybackslash}p{9.0cm}@{}}
  \textbf{Azione} & \textbf{Descrizione} \\
  \textit{Resume}     & riprende il normale flusso di esecuzione, fino al
                        prossimo punto di interruzione. \\
  \textit{Step Over}  & esegue l'istruzione successiva come un blocco unico,
                        senza entrare nei suoi passi interni. \\
  \textit{Step Into}  & entra nell'istruzione successiva per seguirne
                        l'esecuzione riga per riga. \\
  \textit{Step Out}   & essendo dentro un metodo, ne completa le righe restanti
                        e torna al contesto di esecuzione precedente. \\
  \textit{Restart}    & termina l'esecuzione e ricomincia il debug con la
                        configurazione corrente (non funziona con Auto Attach). \\
  \textit{Disconnect} & termina la sessione di debug. \\
\end{tabular}
\end{center}

\info{Meglio del \code{console.log}}{
  È molto comune, all'inizio, cercare di capire un bug disseminando il codice di
  \code{console.log}. Funziona, ma è lento e invasivo. Mezz'ora spesa a imparare
  il debugger si ripaga molte volte: si vedono \textbf{tutte} le variabili in
  scope in un colpo solo, senza modificare il codice e senza dover indovinare in
  anticipo quali stampare.}

\subsection{Ricaricamento automatico}

Ogni volta che modifichiamo il codice dobbiamo riavviare l'intero server
(\code{npm start}, oppure \code{node app.js}). Per rendere lo sviluppo più
piacevole si usano utilità che monitorano il codice e riavviano il server di
sviluppo quando serve: è il \term{live reloading}.

Una di queste è \term{nodemon}:

\begin{lstlisting}[style=shell]
npm install nodemon
\end{lstlisting}

Basta poi sostituire \code{node app.js} con \code{nodemon app.js}:

\begin{lstlisting}[style=json, name=package.json]
"scripts": {
  "start": "nodemon app.js"
},
\end{lstlisting}

\section{Riferimenti}

\begin{itemize}
  \item \textbf{Getting started} --- Node.js Docs. \\
        \urlx{https://nodejs.org/en/learn/} \\
        Parti rilevanti: \textit{Introduction to Node.js}, \textit{How to
        install Node.js}, \textit{Differences between Node.js and the Browser},
        \textit{The V8 JavaScript Engine}, \textit{An introduction to the npm
        package manager}, \textit{ECMAScript 2015 (ES6) and beyond}.
  \item \textbf{Eloquent JavaScript} (3\textsuperscript{a} edizione), Marijn
        Haverbeke, capitolo 20. \\
        \urlx{https://eloquentjavascript.net/}
  \item \textbf{Mixu's Node book}, Mikito Takada. \\
        \urlx{http://book.mixu.net/node/single.html} \\
        Parti rilevanti: capitolo 2 (\textit{What is Node.js?}), capitolo 8
        sezione 8.2 (\textit{NPM}).
  \item \textbf{Nodemon} --- sito ufficiale. \\
        \urlx{https://nodemon.io/}
  \item \textbf{Node.js debugging in VS Code}. \\
        \urlx{https://code.visualstudio.com/docs/nodejs/nodejs-debugging}
\end{itemize}
